%%%  Chapter XI -- Bibliographical Notes (first page)
\section*{BIBLIOGRAPHICAL NOTES}

The following general references may be noted:

\begin{enumerate}
\item \textsc{Sir Horace Lamb}, \emph{Hydrodynamics}, \S\S\,232 and 268,
  Cambridge, England, 1932.

\item J.~\textsc{Proudman}, \emph{Dynamical Oceanography}, chap.\ xii,
  \S\S\,181--3, John Wiley \& Sons, Inc., New York, 1953.
\end{enumerate}

\paragraph{\S\,101.}
The first reference to what has come to be known as the Kelvin--Helmholtz
instability occurs in the writings of Helmholtz:

\begin{enumerate}\setcounter{enumi}{2}
\item H.~\textsc{Helmholtz}, `Ueber discontinuirliche Fl\"ussigkeitsbewegungen',
  \emph{Wissenschaftliche Abhandlungen}, 146--57, J.~A.\ Barth, Leipzig, 1882;
  translation by Guthrie in \emph{Phil.\ Mag.}, Ser.~4, \textbf{36}, 337--46 (1868).
\end{enumerate}

\noindent
(The quotation in the text is from Guthrie's translation.)

Helmholtz's discussion was largely qualitative. The same basic question was
treated analytically by:

\begin{enumerate}\setcounter{enumi}{3}
\item \textsc{Lord Kelvin}, `Hydrokinetic solutions and observations',
  `On the motion of free solids through a liquid, 69--75',
  `Influence of wind and capillarity on waves in water supposed frictionless,
  76--85', \emph{Mathematical and Physical Papers}, iv,
  \emph{Hydrodynamics and General Dynamics}, Cambridge, England, 1910.
\end{enumerate}

\noindent
Kelvin's discussion of the stability of superposed fluids in a state of
differential streaming is exceptionally complete; there is hardly a formula
in this section which cannot be found in Kelvin's paper. The critical wind
velocity, 650~cm/sec (or 660~cm/sec as Kelvin gave it), occurs, for example,
for the first time in this paper. As stated in the text, the significance
of this particular velocity for the generation of waves by wind has been the
subject of much discussion; an account of it will be found in Proudman's book
(reference~2; chap.\ xii). The recognition that Kelvin's critical velocity
may, indeed, have a bearing on phenomena which are observed is more recent:

\begin{enumerate}\setcounter{enumi}{4}
\item W.~H.~\textsc{Munk}, `A critical wind speed for air--sea boundary processes',
  \emph{J.\ Marine Research}, \textbf{6}, 203--18 (1947).
\end{enumerate}

\noindent
For a further discussion of these and related matters, see:

\begin{enumerate}\setcounter{enumi}{5}
\item J.~W.~\textsc{Miles}, `On the generation of surface waves by shear flows',
  \emph{J.\ Fluid Mech.}\ \textbf{3}, 185--204 (1957).

\item ------,
  `On the generation of surface waves by shear flows.~II',
  ibid.\ \textbf{6}, 568--82 (1959).

\item ------,
  `On the generation of surface waves by shear flows.~III.
  Kelvin--Helmholtz instability', ibid.\ 583--98 (1959).
\end{enumerate}

\noindent
The experimental demonstration of the Kelvin--Helmholtz instability
illustrated in the text (Fig.~117) is due to:

\begin{enumerate}\setcounter{enumi}{8}
\item J.~H.~D.~\textsc{Francis}, `Wave motions and the aerodynamic drag on a free oil
  surface', \emph{Phil.\ Mag.}, Ser.~7, \textbf{45}, 695--702 (1954).

\item ------,
  `Wave motions on a free oil surface', ibid., Ser.~8, \textbf{1}, 685--8 (1956).
\end{enumerate}

\paragraph{\S\,102.}
The distributions of $p$ and $U$ considered in this section have been
investigated by:

\begin{enumerate}\setcounter{enumi}{10}
\item \textsc{Sir Geoffrey Taylor}, `Effect of variation in density on the
  stability of superposed streams of fluid',
  \emph{Proc.\ Roy.\ Soc.\ (London)~A}, \textbf{132}, 499--523 (1931).

\item S.~\textsc{Goldstein}, `On the stability of superposed streams of fluids of
  different densities', ibid.\ 524--48 (1931).
\end{enumerate}

\paragraph{\S\,103.}
The case of a heterogeneous fluid with an exponentially decreasing density
and a linearly increasing stream velocity was first considered by Sir~Geoffrey
Taylor (reference~11). The definitive discussion of this case is due to:

\begin{enumerate}\setcounter{enumi}{12}
\item K.~M.~\textsc{Case}, `Stability of an idealized atmosphere.
  I.~Discussion of results',
  \emph{The Physics of Fluids}, \textbf{3}, 149--54 (1960).

\item F.~J.~\textsc{Dyson}, `Stability of an idealized atmosphere.
  II.~Zeros of the confluent hypergeometric function',
  ibid.\ 155--7 (1960).
\end{enumerate}

\noindent
In connexion with Case's investigation (reference~13), see his other papers:

\begin{enumerate}\setcounter{enumi}{14}
\item K.~M.~\textsc{Case}, `Taylor instability of an inverted atmosphere',
  ibid.\ 366--8 (1960).

\item ------,
  `Stability of inviscid plane Couette flow', ibid.\ 143--8 (1960).
\end{enumerate}

\paragraph{\S\,104.}
The example treated in this section is due to:

\begin{enumerate}\setcounter{enumi}{16}
\item P.~G.~\textsc{Drazin}, `The stability of a shear layer in an unbounded
  heterogeneous inviscid fluid',
  \emph{J.\ Fluid Mech.}\ \textbf{4}, 214--24 (1958).
\end{enumerate}

\paragraph{\S\,105.} See:

\begin{enumerate}\setcounter{enumi}{17}
\item S.~\textsc{Chandrasekhar}, `The effect of rotation on the development of
  the Kelvin--Helmholtz instability', unpublished.
\end{enumerate}

\noindent
For a different treatment of the same problem see:

\begin{enumerate}\setcounter{enumi}{18}
\item H.~\textsc{Solberg}, `Integrationen der atmosph\"arischen
  St\"orungsgleichungen, I.
  Wellenbewegungen in rotierenden, inkompressiblen Fl\"ussigkeitsschichten',
  \emph{Geofysiske Publikasjoner}, \textbf{9}, 1--120 (1928).
\end{enumerate}

\noindent
See also:

\begin{enumerate}\setcounter{enumi}{19}
\item V.~\textsc{Bjerknes}, J.~\textsc{Bjerknes}, H.~\textsc{Solberg},
  and T.~\textsc{Bergeron},
  \emph{Physikalische Hydrodynamik}, chap.~11, Springer, Berlin, 1933.
\end{enumerate}

\paragraph{\S\,106.}
The Kelvin--Helmholtz instability in the framework of hydromagnetics has
been considered by:

\begin{enumerate}\setcounter{enumi}{20}
\item D.~H.~\textsc{Michael}, `The stability of a combined current and vortex sheet
  in a perfectly conducting fluid',
  \emph{Proc.\ Camb.\ Phil.\ Soc.}\ \textbf{51}, 528--32 (1955).
\end{enumerate}
